% main.tex
%\documentclass{gtac}
\documentclass[blind]{gtac}

\usepackage{cite}
\usepackage{xcolor}
\usepackage{booktabs}
\usepackage{multirow}
\usepackage{makecell}
\usepackage{graphicx}
\usepackage{amsmath}
\usepackage{tikz}
\usetikzlibrary{positioning, shapes.geometric, arrows.meta}
\usepackage[colorlinks=true, linkcolor=blue, urlcolor=blue, citecolor=blue]{hyperref}
\renewcommand\theadfont{\bfseries\small}
\renewcommand\theadgape{}

\SetWatermarkScale{0.3}
\SetWatermarkAngle{30}
\SetWatermarkColor{gray!20}

\title{Orbit: MPC-Augmented Routing for Heterogeneous LLM Serving --- Design, Evaluation, and Lessons Learned\subtitle{Submitted Anonymously for Double-Blind Review}}

\author{Author 1, Author 2, Author 3 \\
AMD \\
email1@amd.com, email2@amd.com, email3@amd.com}

\begin{document}

\maketitle

\pagestyle{fancy}
\fancyhf{}
\cfoot{\thepage}
\renewcommand{\headrulewidth}{0pt}
\chead{{\sffamily\textcolor{red}{\fontsize{10}{12}\selectfont\textbf{DO NOT FORWARD – FOR GTAC'26 REVIEW COMMITTEE ONLY}}}}
\thispagestyle{fancy}

\begin{abstract}
Heterogeneous LLM inference clusters — where servers differ in GPU generation,
thermal state, or concurrency capacity — expose a gap in existing routers that
treat all servers as identical. We design and implement \textbf{Orbit}, a
request router that augments Power-of-Two-Choices (PO2) with a Model Predictive
Control (MPC) layer that estimates per-server service rates and computes
capacity-aware routing weights via a short-horizon quadratic program.
We evaluate Orbit on AMD MI300X GPUs with five policies (Round Robin, Least
Outstanding Requests, PO2, MPC-PO2, and a new Kalman-PO2) across a
1$\times$--8$\times$ heterogeneity sweep and a 10,000-request convergence study.
Our central finding is a \emph{negative result}: both MPC-PO2 and Kalman-PO2
are consistently 15--40\% worse than PO2 at steady state across all tested
conditions, and the gap \emph{widens} rather than closes under overload
(4--12\,rps). Kalman-PO2 shows a transient advantage in short runs
(327\,ms vs.\ 386\,ms mean, $10\times$ lower P99) by eliminating PO2's
burst variance, but this benefit does not persist beyond $\approx$500
requests. An overload study (up to 12\,rps, 3$\times$ baseline) falsifies
the hypothesis that PO2 breaks down under load: LOR and PO2 remain the
most robust policies throughout.
We identify the root cause --- both MPC and Kalman measure
\emph{experienced} load rather than \emph{intrinsic} capacity, leaving no
exploitable signal beyond what PO2 already captures --- and derive
concrete redesign directions. We present this as a cautionary case study
in applying control-theoretic methods to LLM serving.
\end{abstract}

\keywords{LLM serving, request routing, load balancing, model predictive control,
Kalman filter, heterogeneous inference, vLLM, ROCm, AMD MI300X, negative result}

% -----------------------------------------------------------------------
\section{Introduction}
% -----------------------------------------------------------------------

GPU clusters serving LLMs exhibit significant heterogeneity in practice:
servers run on different GPU generations, experience variable thermal
throttling, hold different amounts of cached KV data, and operate under
mixed tenant workloads. A router that ignores these differences will
overload slower servers, inflating tail latency even when aggregate capacity
is sufficient.

Reactive policies such as Least Outstanding Requests (LOR) and
Power-of-Two-Choices (PO2) \cite{po2} respond to observed queue depth
instantaneously but cannot, in principle, distinguish a server that is
slow because it is overloaded from one that is intrinsically lower capacity.
This motivates a predictive approach: if a router can \emph{learn} each
server's throughput, it can steer traffic proportionally to capacity rather
than reactively to queue depth.

We present \textbf{Orbit}, which implements this idea by wrapping PO2 with
a Model Predictive Control (MPC) \cite{mpc_book} layer. Orbit estimates
per-server service rates via exponential moving averages (EMA) and solves a
quadratic program (QP) every 100\,ms to compute routing weights that drive
queue depths toward a target. The implementation integrates with AMD's vLLM
ROCm stack and supports both standalone and prefill-decode disaggregated
topologies.

We evaluate Orbit thoroughly and report an \textbf{honest negative result}:
MPC-PO2 does not outperform PO2 or LOR under any tested condition, and
the performance gap is persistent --- it does not close with more traffic.
We believe this result is more valuable than a selective positive report
because it exposes \emph{why} queue-depth-based MPC fails for LLM routing
and what signals would actually work.

% -----------------------------------------------------------------------
\section{Design}
% -----------------------------------------------------------------------

\subsection{System Architecture}

Orbit is a FastAPI proxy between clients and a pool of vLLM servers.
All experiments in this paper use \emph{single-pool} mode: each request
is dispatched to one server that handles both prefill and decode.
The router tracks per-server inflight counts and applies one of five
policies: \textbf{RR} (round-robin), \textbf{LOR} (min inflight),
\textbf{PO2} (sample two, pick lower weighted-inflight score),
\textbf{MPC-PO2} (PO2 with MPC-computed weights), and
\textbf{Kalman-PO2} (PO2 with Kalman-filtered latency scores).
Orbit also supports prefill-decode disaggregated topologies as a
deployment option, though evaluation of that mode is left for future work.

\subsection{MPC Weight Controller}

Every $\Delta t = 100$\,ms, Orbit solves a per-server QP to update weights
$w_i$. Let $\hat{\mu}_i$ be the EMA service rate ($\alpha = 0.3$),
$\bar{q}$ the mean inflight count, and $q^* = 1.0$ the normalized target.
The horizon-$H$ QP is:

\begin{align}
\min_{w,\, q} \quad & \sum_{k=0}^{H-1} \!\left[ (q_{k+1} - q^*)^2
    + \lambda(w_k - 1)^2 \right] \label{eq:mpc} \\
\text{s.t.} \quad
  & q_{k+1} = q_k + \Delta t \cdot
    \tfrac{\lambda_a p_i - \hat{\mu}_i}{\bar{q}},\notag\\
  & q_{k+1} \geq 0,\quad w_{\min} \leq w_k \leq w_{\max} \notag
\end{align}

\noindent where $\lambda_a$ is the arrival rate, $p_i = w_k/\sum_j w_j$
is server $i$'s traffic fraction, and $\lambda=0.07$ penalizes weight
deviation. We use $H=10$, $w_{\min}=0.1$, $w_{\max}=3.0$, solved with
OSQP \cite{osqp} in under 1\,ms. The regularization $\lambda(w_k-1)^2$
prevents extreme weights during estimation warm-up \cite{mpc_book}.

% -----------------------------------------------------------------------
\section{Evaluation}
\label{sec:results}
% -----------------------------------------------------------------------

\subsection{Experimental Setup}

All experiments run on AMD MI300X nodes (192\,GB HBM3/GPU) with
vLLM v0.23.0 ROCm, serving Qwen3-8B in bfloat16 (\texttt{max-model-len=2048},
\texttt{gpu-memory-utilization=0.15}). Heterogeneity is simulated via two vLLM servers on separate GPUs
by varying \texttt{max-num-seqs} --- vLLM's concurrency cap on how
many sequences are batched per iteration.
\textbf{Server 0:} \texttt{max-num-seqs=64} (fixed);
\textbf{Server 1:} \texttt{max-num-seqs} $\in \{64,32,16,8\}$,
yielding 1$\times$--8$\times$ capacity ratios.
Importantly, this creates a \emph{queue-depth} bottleneck rather than
a per-request compute speed difference: at light load (4\,rps) both
servers process individual requests in $\approx$280\,ms, but the
constrained server's queue fills faster under burst arrivals, producing
occasional high-latency outliers that inflate P99.
Requests follow Poisson arrivals at 4\,req/s with 15 fixed prompts
(\texttt{max\_tokens=50}). The router is freshly restarted before each
policy run so MPC starts with no prior EMA state.

\subsection{Heterogeneity Sweep}

Table~\ref{tab:setA} shows results for 120 requests across all four
heterogeneity levels.
\textbf{RR} shows high P99 (1775--2587\,ms) from periodically routing
to the slow server regardless of capacity.
\textbf{LOR and PO2 perform comparably} across all ratios
(mean 285--294\,ms), with PO2's random sampling introducing minor noise.
\textbf{MPC-PO2 is 14--20\% worse} than PO2 at every heterogeneity
level, with elevated stdev (42--50\,ms vs.\ 17--20\,ms for PO2).
Critically, the MPC overhead is \emph{independent} of the heterogeneity
ratio --- the same penalty at 1$\times$ as at 8$\times$ --- suggesting
the problem is not the degree of heterogeneity but the nature of the
routing signal.

\begin{table}[h]
\centering
\caption{Heterogeneity sweep: 4\,rps, 120 req, fresh router per policy.}
\label{tab:setA}
\small
\begin{tabular}{@{}c l r r r@{}}
\toprule
\thead{Hetero} & \thead{Policy} & \thead{Mean\\(ms)} & \thead{Stdev\\(ms)} & \thead{P99\\(ms)} \\
\midrule
\multirow{4}{*}{1$\times$}
  & RR      & 413.8 & 431.7 & 2587.1 \\
  & LOR     & \textbf{292.3} & 19.7 & 352.9 \\
  & PO2     & 289.0 & 17.8 & 323.5 \\
  & MPC-PO2 & 341.8 & 47.4 & 444.2 \\
\midrule
\multirow{4}{*}{4$\times$}
  & RR      & 354.0 & 313.7 & 1828.4 \\
  & LOR     & \textbf{285.2} & 17.7 & 330.9 \\
  & PO2     & 285.9 & 19.7 & 344.4 \\
  & MPC-PO2 & 326.4 & 45.8 & 429.4 \\
\midrule
\multirow{4}{*}{8$\times$}
  & RR      & 353.8 & 338.1 & 2417.1 \\
  & LOR     & \textbf{287.9} & 16.7 & 330.8 \\
  & PO2     & 287.9 & 17.8 & 342.4 \\
  & MPC-PO2 & 346.6 & 49.6 & 448.4 \\
\bottomrule
\end{tabular}
\end{table}

\subsection{Convergence Study: 10,000 Requests}

To test whether MPC eventually converges past a cold-start period,
we ran PO2, MPC-PO2, \emph{and} Kalman-PO2 for \textbf{10,000 requests}
at 4\,rps with 8$\times$ heterogeneity ($\approx$42 minutes each),
using fresh vLLM containers per policy.
Table~\ref{tab:conv} shows per-200-request windows.
The MPC gap \emph{never closes}: MPC-PO2 is 13--22\% worse than PO2
from the first window to the last, with no downward trend.
Kalman-PO2 similarly settles at 15--21\% above PO2 throughout,
confirming that the advantage seen in the 500-request pilot
(\S\ref{sec:kalman}) does not persist to steady state.

\begin{table}[h]
\centering
\caption{Convergence study (8$\times$ hetero, 4\,rps, 10\,K req).
  Selected windows of 200 requests; full 50-window table in
  accompanying repository.}
\label{tab:conv}
\small
\begin{tabular}{@{}r r r r r@{}}
\toprule
\thead{Req\\window} & \thead{PO2\\(ms)} & \thead{MPC\\(ms)} & \thead{Kalman\\(ms)} & \thead{MPC vs\\PO2} \\
\midrule
0--200     & 537.9 & 418.8 & 612.4 & \textbf{+22.1\%} \\
200--400   & 289.7 & 330.7 & 325.3 & $-$14.2\% \\
1000--1200 & 288.2 & 336.7 & 332.7 & $-$16.8\% \\
3000--3200 & 289.3 & 331.7 & 331.9 & $-$14.7\% \\
5000--5200 & 288.4 & 332.5 & 342.7 & $-$15.3\% \\
7000--7200 & 282.9 & 330.5 & 336.3 & $-$16.8\% \\
9800--10000& 282.9 & 335.6 & 324.4 & $-$18.6\% \\
\midrule
\textbf{Overall} & \textbf{291.8} & \textbf{338.2} & \textbf{342.9}
  & $-$\textbf{15.9\%} \\
\bottomrule
\end{tabular}
\end{table}

\subsection{Root Cause Analysis}

The convergence failure reveals a fundamental mismatch between the MPC
design and the routing problem. At 4\,rps with 8$\times$ heterogeneity,
our log analysis shows PO2 routes 49\%/51\% across the two servers by
request count, yet achieves mean latency of 289\,ms --- almost identical
to LOR. This is because PO2 \emph{routes opportunistically}, sending
short requests to whichever server happens to be available.  The slow
server's higher latency means it is sampled as higher-load more often,
so PO2 naturally deflects the longest-running requests away from it.
The EMA rate estimator, operating on a 300\,ms window at 4\,rps, sees
fewer than 1 completion per server per window --- insufficient to
distinguish a slow server from a temporarily busy fast server.  There is
effectively no exploitable signal for MPC.

MPC's overhead then makes things worse: the regularization term
$\lambda(w_k-1)^2$ biases weights toward 1.0, occasionally sending
extra traffic to the slow server that PO2 would have deflected.

\subsection{Overload Study: 4--12\,rps}
\label{sec:overload}

One might expect that PO2's implicit balancing breaks down under high
load, where queues saturate and queue-depth ceases to be informative.
We test this directly by sweeping arrival rates from 4 to 12\,rps
(3$\times$ above baseline) at 8$\times$ heterogeneity, 500 requests
per cell, fresh containers per policy. Table~\ref{tab:overload} shows
the results.

\begin{table}[h]
\centering
\caption{Overload sweep: 8$\times$ hetero, 500\,req per cell.
  vs\_PO2 column shows MPC-PO2 and Kalman-PO2 mean relative to PO2.}
\label{tab:overload}
\small
\begin{tabular}{@{}c l r r r r@{}}
\toprule
\thead{Rate} & \thead{Policy} & \thead{Mean\\(ms)} & \thead{Stdev\\(ms)} & \thead{P99\\(ms)} & \thead{vs\\PO2} \\
\midrule
\multirow{4}{*}{4\,rps}
  & LOR     & \textbf{314.4} & 203.7 & 1620 & --- \\
  & PO2     & 379.3 & 560.9 & 4688 & --- \\
  & MPC-PO2 & 360.9 & 239.9 & 2082 & $-$4.8\% \\
  & Kalman  & 356.0 & 223.4 & 1965 & $-$6.1\% \\
\midrule
\multirow{4}{*}{6\,rps}
  & LOR     & 384.9 & 538.5 & 3766 & --- \\
  & PO2     & \textbf{326.5} & 224.1 & 1914 & --- \\
  & MPC-PO2 & 450.5 & 586.8 & 4247 & +38.0\% \\
  & Kalman  & 380.1 & 247.7 & 2042 & +16.4\% \\
\midrule
\multirow{4}{*}{8\,rps}
  & LOR     & \textbf{416.5} & 574.1 & 4314 & --- \\
  & PO2     & 445.9 & 642.2 & 4675 & --- \\
  & MPC-PO2 & 479.3 & 600.4 & 4193 & +7.5\% \\
  & Kalman  & 500.9 & 632.1 & 4492 & +12.3\% \\
\midrule
\multirow{4}{*}{12\,rps}
  & LOR     & \textbf{456.6} & 637.7 & 4892 & --- \\
  & PO2     & 463.5 & 609.3 & 4345 & --- \\
  & MPC-PO2 & 541.6 & 704.1 & 4602 & +16.8\% \\
  & Kalman  & 643.5 & 755.6 & 3923 & +38.8\% \\
\bottomrule
\end{tabular}
\end{table}

The overload hypothesis is \textbf{falsified}: MPC-PO2 and Kalman-PO2
do not recover under load --- they get \emph{worse}. At 12\,rps,
Kalman-PO2 is 39\% above PO2 in mean latency.
A 2,000-request convergence run at 12\,rps confirms the gap is
persistent: PO2=376\,ms, MPC-PO2=421\,ms (+12\%), Kalman-PO2=485\,ms
(+29\%), with no downward trend across ten 200-request windows.

Two additional observations: (1)~\textbf{LOR emerges as the most
robust policy under overload} --- it ties or beats PO2 at 8--12\,rps
while MPC/Kalman degrade. LOR's per-request inflight count is a more
reliable signal than either EMA rate estimates or filtered latency when
queues are deep. (2)~At 4\,rps baseline this run, PO2 exhibits its
known burst variance (P99=4,688\,ms), giving MPC and Kalman a
spurious short-run advantage --- further evidence that the short-run
500-request result (\S\ref{sec:kalman-short}) should not be
over-interpreted.

% -----------------------------------------------------------------------
\section{Related Work}
% -----------------------------------------------------------------------

\textbf{LLM Serving.}
PagedAttention \cite{vllm} and continuous batching \cite{orca} established
the modern vLLM serving stack. Disaggregated PD architectures
\cite{splitwise,distserve,mooncake} separate prefill and decode for
independent scaling. Sarathi-Serve \cite{sarathi} addresses prefill-decode
interference via chunked prefills.

\textbf{Load Balancing.}
PO2 \cite{po2} achieves $O(\log \log n)$ maximum queue depth.
vLLM-router \cite{vllm_router} implements round-robin, PO2, consistent
hashing, and cache-aware routing with an approximate radix tree for prefix
affinity. Our work shows that naive MPC over queue-depth signals does not
improve on PO2, and identifies what richer signals would be needed.

\textbf{MPC for Systems.}
MPC \cite{mpc_book} with OSQP \cite{osqp} enables sub-millisecond
constrained optimization. Its application to LLM routing is nascent;
our study characterizes the conditions under which a queue-depth feedback
signal is insufficient.

% -----------------------------------------------------------------------
\section{Kalman-PO2: Validating the Diagnosis}
\label{sec:kalman}
% -----------------------------------------------------------------------

The root cause analysis points to a concrete fix: replace queue-depth
(which is masked by PO2's implicit balancing) with \emph{observed
request latency} as the routing signal. Latency remains discriminative
even when queues are balanced, because a capacity-constrained server
(\texttt{max-num-seqs}$=8$) takes longer per request regardless of
queue depth --- it batches fewer requests per step.

We implement \textbf{Kalman-PO2}: a PO2-style router where the
server-selection score is the Kalman-filtered estimate of each server's
mean request latency, rather than weighted inflight count. The Kalman
state per server is:
\[
  \hat{L}_{k|k} = \hat{L}_{k-1|k-1} + K_k(z_k - \hat{L}_{k-1|k-1}),
  \quad K_k = \frac{P_{k|k-1}}{P_{k|k-1} + R}
\]
where $z_k$ is the observed latency of completed request $k$,
$Q=400\,\text{ms}^2$ (process noise), $R=3600\,\text{ms}^2$
(measurement noise), $\hat{L}_0=350\,\text{ms}$, and
$P_0=40{,}000\,\text{ms}^2$ (high initial uncertainty for fast cold
start). Routing: sample two servers, pick lower $\hat{L}$. No QP, no
arrival rate estimation, no regularization.

\subsection{Short-Run Results (500 Requests)}
\label{sec:kalman-short}

Table~\ref{tab:kalman} compares all three policies at 8$\times$
heterogeneity with 500 warm-server requests. Kalman-PO2 achieves lower
mean (327\,ms vs.\ 386\,ms) and dramatically lower P99 (410\,ms vs.\
4,527\,ms for PO2), and outperforms MPC-PO2 as well.

\begin{table}[h]
\centering
\caption{8$\times$ heterogeneity, 4\,rps, 500\,req. Kalman filter
         converged estimates: fast=309\,ms, slow=366\,ms.}
\label{tab:kalman}
\small
\begin{tabular}{@{}l r r r@{}}
\toprule
\thead{Policy} & \thead{Mean\\(ms)} & \thead{Stdev\\(ms)} & \thead{P99\\(ms)} \\
\midrule
PO2        & 386.6 & 583.9 & 4527.2 \\
MPC-PO2    & 334.2 &  44.2 &  446.2 \\
\textbf{Kalman-PO2} & \textbf{327.0} & \textbf{36.3} & \textbf{410.1} \\
\bottomrule
\end{tabular}
\end{table}

Log-level analysis reveals the mechanism and its limitations. With
only two servers, \texttt{random.sample(pool,\,2)} always returns
both, so routing is deterministic: always pick the server with the
lower Kalman estimate. Three phases emerge: (i)~\emph{convergence}
(req~0--50, 80\% fast), (ii)~\emph{stable fast-routing}
(req~50--300, 100\% fast, mean$\approx$309\,ms), and
(iii)~\emph{partial inversion} (req~300--500, $\sim$64\% fast,
mean$\approx$330\,ms). During phase~(iii), concentrating load on
the fast server inflates \emph{its} observed latency above the slow
server's, causing the Kalman estimate to partially invert. This
feedback instability is a fundamental limitation: a latency signal
measures \emph{experienced} latency, not \emph{intrinsic} capacity.
Kalman-PO2 nonetheless wins over PO2 in this run because PO2's
variance is much larger ($\sigma=584$\,ms): PO2 produces 16 requests
$>$1\,s (4.4\%) and a 4,527\,ms P99 from periodic bursts to the slow
server.

\subsection{Long-Run Behavior (10,000 Requests)}

The 10,000-request convergence study (\S\ref{sec:results}, Table~\ref{tab:conv})
reveals the steady-state picture. With a fresh server restart,
Kalman-PO2 settles into an equilibrium routing $\approx$62\% of
requests to the fast server. This stable partial-routing avoids the
burst inversions seen in the short run, but it also eliminates
Kalman's advantage: at steady state, mean=342.9\,ms vs.\
PO2=291.8\,ms ($-$17.5\%). The gap is indistinguishable from MPC-PO2
(338.2\,ms). Both Kalman and MPC remain 15--20\% worse than PO2 for
the entire 10,000-request run, with no convergence.

The explanation is the same as for MPC: Kalman's latency signal reflects
\emph{queue state}, not capacity. Under steady load at 4\,rps,
the fast server's estimate rises as it absorbs more traffic, and
the system finds a sub-optimal equilibrium rather than converging
to a true capacity-proportional split.

On homogeneous clusters (1$\times$), Kalman-PO2 incurs a 6.1\%
overhead (315\,ms vs.\ 297\,ms) from the same concentration effect.
An uncertainty-gated fallback to PO2 when
$|\hat{L}_a - \hat{L}_b| < \sqrt{P_a + P_b}$ would address both
pathologies --- left for future work.

% -----------------------------------------------------------------------
\section{Conclusions}
\label{sec:conclusions}
% -----------------------------------------------------------------------

We show through experiment and repeated replication that the
\emph{choice of feedback signal} is the critical variable in LLM
request routing. Queue-depth is a sufficient statistic when load is
moderate and servers are stationary --- PO2 already exploits it
optimally, leaving nothing for MPC to improve.

Switching to latency as the feedback signal (Kalman-PO2) yields a
transient benefit: in short runs ($\leq$500 requests), Kalman-PO2
eliminates PO2's burst variance ($\sigma$: 584\,ms$\to$36\,ms) and
achieves a $10\times$ lower P99 (4,527$\to$410\,ms). But the
10,000-request replication reveals that this advantage does not
persist: Kalman-PO2 settles at the same $\sim$16\% overhead as MPC-PO2.
Both signals measure \emph{experienced} rather than \emph{intrinsic}
latency, and under sustained load PO2's implicit balancing eliminates
any exploitable differential.

\textbf{Takeaways for practitioners:}
(1) \textbf{PO2 and LOR are robust across all load regimes tested} ---
from moderate (4\,rps) to 3$\times$ overload (12\,rps). Adding MPC or
Kalman complexity makes things worse, not better, at every scale.
(2) \textbf{LOR is underrated at high load}: it ties or beats PO2 at
8--12\,rps where PO2's random sampling occasionally misdirects requests.
(3) \textbf{Signal fidelity matters more than controller sophistication}:
both MPC's EMA rate estimates and Kalman's filtered latency measure
load state rather than intrinsic capacity, making them systematically
exploitable by the very routing decisions they drive.
(4) A feedback signal that encodes intrinsic capacity directly
(e.g.\ time-to-first-token, or throughput under a controlled probe)
is the missing ingredient for any predictive router to outperform PO2.

\textbf{Open directions:} uncertainty-gated fallback to PO2 when
Kalman estimates overlap within their posterior variance; TTFT-specific
routing for prefill-decode disaggregated topologies where queue depth
and latency decouple; and probing-based capacity estimation to break
the fundamental signal limitation identified here.

\section*{\st{Acknowledgments} Do Not Include in Submission}
Do \textcolor{red}{\ul{\textbf{not}}} include acknowledgments in this blind-review version.

\begingroup
\let\itshape\upshape
\bibliographystyle{abbrv}
\bibliography{refs}
\endgroup

\end{document}
